% ===== 01_introduction =====
\section{引言}
\label{sec:introduction}
限制感染高峰是医疗资源有限时的一项防控目标，但实现同一峰值上限的干预组合可能具有不同的控制时长、累计感染和隔离负担。Ferguson 等\cite{Ferguson2020Report9}在 COVID-19 早期的建模研究中，将降低医疗需求高峰作为缓解策略的目标，并指出需要多种措施配合。严格措施的持续时间与解除后的传播同样影响政策选择，Duque 等\cite{Duque2020}据此研究了在避免医院超载的同时缩短干预时间的触发规则。因此，感染上限需要与干预的启动、分配和退出要求一同评价。

容量约束控制为这一问题提供了数学形式。Miclo 等\cite{Miclo2022ICU}在 SIR 模型、ICU 上限和特定线性累计成本下求得解析最优策略：先按常规水平传播，首次触及上限后维持恒定感染，逐步放松措施，直至恢复常规水平。Avram 等\cite{Avram2022}分析了有限干预能力下的可行区域及退出目标，Angulo 等\cite{Angulo2021}给出峰值约束下的干预能力与时长判据，能力不足时可能需要提前启动。Balderrama 等\cite{Balderrama2024}在累计干预和感染峰值限制下最小化最终感染规模，其策略也包含恒定感染阶段。恒定感染因而是容量约束控制中的一种已有结构，其最优性取决于模型、成本和退出任务。

接触减少与追踪隔离通过不同机制降低传播。前者降低接触频率，后者识别病例及其接触者，并将被追踪者移出传播人群。追踪效果取决于覆盖比例和响应速度\cite{Hellewell2020,Ferretti2020}。Kucharski 等\cite{Kucharski2020}比较多种措施组合时，同时估计需要隔离的接触者数量，使传播效果与隔离负担得到分别评价。Tang 等\cite{Tang2020Transmission}在含接触追踪的 SEIR 型模型中，将被追踪的感染者与未感染者分入不同仓室。这一区分使未感染者隔离成为独立的计量对象。

Zhou 等\cite{Zhou2024SPCI}在含追踪隔离的 SIR 型模型中采用分段常数接触率和隔离比例，以首次积分分析干预时机、感染峰值和隔离区外清零时间。其模型区分非隔离与隔离感染者，社区传播由未被隔离的感染者引起。He 等\cite{He2023TDINN}将动力学与神经网络结合，从包括西安在内的疫情数据重构随时间变化的接触率和隔离比例。这些工作分别提供了分段干预的解析方法与时间控制函数的数据参照。本文模型中进入 $I$ 和 $I_q$ 的全部新增感染均由社区中的 $I$ 引起；接触率有界时，限制 $I$ 也为总新增感染流量提供上界。

多种干预与容量目标的结合也已有研究。Charpentier 等\cite{Charpentier2020}在区分检出与未检出人群的扩展模型中纳入 ICU 容量，数值研究检测与封控的分配。Pollinger\cite{Pollinger2023}分析接触追踪与社交距离的联合抑制，联系追踪效率、干预强度和持续时间。Balakrishnan 和 Palathingal\cite{Balakrishnan2023}通过动态调整检测率维持给定医疗资源使用目标。本文关注其中一个具体的分配问题：在区分感染与未感染者隔离的模型中，刻画同一进入状态、社区感染上限和退出状态下的全部阈值控制。单一传播系数将两种措施合并，分段控制和数据重构则选取特定路径；分别描述两种措施并考察全部允许分配，才能比较未感染者隔离与控制时长、累计感染及成本之间的关系。

本文采用隔离易感者在研究时段内不返回社区的 SIQR 型模型。记社区易感者与非隔离感染者人数为 $S(t),I(t)$，接触率与追踪隔离比例为 $c(t),q(t)$。社区感染上限 $\eta$ 是外生设计参数，医疗容量是选取它的动机。维持 $I(t)\equiv\eta>0$ 要求
\[
\frac{\beta c(t)(1-q(t))S(t)}{N}=\gamma,
\]
其中 $\beta$ 为传播概率，$\gamma$ 为非隔离感染者移出率，$N$ 为总人口。固定 $c=c_0$ 时，该关系唯一确定 $q$；两者同时可调时，它在每个 $S$ 上只约束乘积 $c(1-q)$。改变分配还会改变未感染者进入隔离的流量，进而改变社区易感者下降的速度与退出时刻。因此，同一峰值约束并不能唯一确定全过程的代价。

在常规轨迹满足触发条件\eqref{eq:model:trigger}时，本文固定首次触及阈值的二维状态 $(S^*,\eta)$，要求随后维持 $I=\eta$，并在首次达到常规传播临界状态 $(S_c,\eta)$ 后恢复常规控制。第\ref{sec:model:control-class}节将这些有界可测控制记为 $\mathcal A(\eta)$。需要回答的问题是：能力限制下能否完成整段控制，不同分配如何改变控制时长、累计新增感染与新增隔离易感者人数，以及二次成本最低的分配是否也最快或感染最少。

仅隔离控制给出这一比较的解析基准。本文由首次积分和 Lambert $W$ 函数确定启动状态，连接阈值控制轨迹、退出后的动态与累计新增感染（第\ref{sec:s1}节）。在触发区间内，提高阈值降低最大隔离比例、控制时长和成本；当 $\eta>1$ 时，计算至清零终点的累计新增感染随之增加。利用阈值控制期内 $S$ 严格下降的性质，联合控制可全部表示为接触率状态函数 $c_c(S)$，能力界转化为逐状态允许区间，由此得到整段可行的充要条件（定理\ref{thm:joint:representation}与推论\ref{cor:joint:feasibility}）。显式策略族用于连接两种单控制，指标比较与成本最小化仍考察整个允许控制类。

固定二维起止状态且 $0<\beta<1$ 时，控制期累计新增感染与控制时长同序，新增隔离易感者人数反序（定理\ref{thm:joint:comparison}）。无额外能力限制时，仅隔离控制最快且控制期累计新增感染最少，同时使最多未感染者进入隔离。同一完整初值、阈值 $\eta>1$ 与前后常规控制下，时长差还决定清零时间差和全过程累计新增感染差（推论\ref{cor:joint:whole-epidemic}）。这些恒等式把控制时长与感染、隔离人数的比较联系起来。

对于正权重的二次强度成本，在不含控制变化率或额外累计预算等跨状态约束的允许控制类中，本文证明最小值可达到，并将其确定归结为逐状态的端点与有限驻点比较（定理\ref{thm:joint:cost-minimum}）。固定其余条件，在触发且整段可行的阈值范围内，最小成本随阈值严格下降。无额外能力限制时，最小成本分配在退出附近共同使用两种措施，其成本严格低于两种单控制各自的成本（命题\ref{prop:joint:terminal-allocation}）。规定有限时刻达到 $S=S_c$ 还给出从初始时刻到退出的累计新增感染下界 $\beta(S_0-S_c)$（命题\ref{prop:exit-infection-lower-bound}）。因此，峰值上限、感染总数和最低成本分别对应不同评价，退出目标也会限制可达到的感染结局。

数值部分在基准与西安参数下，以相同完整初值、阈值和启动、退出要求比较四种分配。已有计算展示二次成本与控制时长的取舍，并分别保留有限能力代表点、有限权重网格和离散匹配乘子点的证据范围。西安 TDINN 重构作为外部参照，其峰值与退出任务不同；设计阈值下的长时结果用于解释规定退出目标的代价。

第\ref{sec:model}节定义模型、控制类与指标，第\ref{sec:s1}节给出仅隔离解析基准，第\ref{sec:joint}节分析联合可行性与指标权衡，第\ref{sec:joint:cost}节确定二次成本分配。第\ref{sec:numerics}节给出同条件数值比较，第\ref{sec:xian}节讨论西安情景与退出代价，第\ref{sec:discussion}节解释适用范围并给出结论。完整证明和辅助分析分别列于附录与补充材料。
